% Generated by tools/textbook/export_listings.js from tools/concepts/assessments.json: do not edit.
\begin{chapstatus}{Il filo del capitolo}
Un registro è fatto di flip-flop che campionano su un fronte del clock: è così che il contatore della Lezione 02 sopravviveva da un frame all'altro, in hardware. Prossimo passo: come i registri si scambiano i dati.
\par\medskip\textbf{Biglietto d'uscita} (da due a quattro minuti, da soli, su carta):
\begin{enumerate}[nosep,leftmargin=2em]
\item Il clock del kit va a 160 MHz. Qual è il suo periodo, e perché ogni rete combinatoria tra due registri deve stabilizzarsi entro quel tempo?
\item Qual è la differenza tra un latch D e un flip-flop D, e quale dei due permette di leggere e scrivere un registro nello stesso ciclo di clock in sicurezza?
\end{enumerate}
\end{chapstatus}
